\documentclass[11pt]{article}
\usepackage[T1]{fontenc}
\usepackage[margin=1in]{geometry}
\usepackage[hidelinks]{hyperref}

\usepackage{amsthm}
\setlength{\parskip}{0.6em}
\setlength{\parindent}{0em}

\title{Omega: A Programming Language Built for Verified Programs}
\author{M. Drake Stapleton}
\date{October 2026}

\begin{document}
\maketitle

\begin{abstract}
We propose Omega, a planned systems programming language that combines
the memory safety of Rust with the simplicity of Go, and
adds two properties neither language has: it is semantics-first, and
verification is a first-class language feature. In Omega, the primary
object is meaning, not text. Syntax is a disposable skin over a canonical,
content-addressed semantic identity: two surface texts that elaborate into
the identical normalized semantic graph share one semantic identity
(SEMANTIC\_ID). The module system is the crumb store. A \emph{crumb} is
a verified program together with its verification receipt, and programs
may only be built from verified crumbs. The doctrine that the verified
program is the universal artifact of meaning is therefore not a policy
document. It is the grammar of the language. In this document we state the
motivation, the design principles, the intended trade-offs, and the open
questions, and then we answer the objections we expect. It is a design
proposal, not a specification. Nothing here is final until we have built,
measured, and ratified it.
\end{abstract}

\section{Motivation}

We propose a programming language built for verified programs.

Most programming languages were designed for a world in which humans write
all the code. That world is ending. In our work, the large majority
of code is already written by AI agents, reviewed by AI agents, and tested by
AI agents. Languages optimized for human comfort are optimizing for the
wrong writer.

Rust is the closest existing answer to safe systems programming, but its
compiler is a harsh teacher: the learning curve is steep, builds are slow,
and the strictness that makes it safe also makes it unfriendly. Go chose the
opposite trade: simple, fast to build, friendly. It keeps memory safety
apart from data races, and pays for simplicity with a garbage collector
and less fine-grained control over layout and timing.

Omega starts from a different question, and we ask it directly: what would
a language look like if we designed it for verifiability first, with the
knowledge that its primary authors will be machines? A strict compiler is
painful for a human.
For an agent, a strict compiler is a tireless reviewer that never sleeps,
never gets tired, and catches the bug before it ships. The pain of
strictness is almost entirely human pain. We design Omega to spend that
budget where it buys the most: correctness you can check, not comfort you
can feel.

The name is deliberate. In our architecture, Omega is the umbrella system
that defines, synthesizes, and coordinates verification. We give the
language the same name because the language is one interface of that
system. Verification itself is performed by an independent verifier,
never by the generator: the language does not grade its own work.

\section{Semantics first}
\label{sec:semantics}

In every mainstream language, the primary object is text. The compiler
reads characters, and meaning is something the text is hoped to carry.
Omega inverts this: the primary object is meaning, and text is a
disposable skin.

This is not a metaphor. It is the founding position of the Omega engine,
established in Milestone 4 (OMEGA\_SEMANTICS): Omega defines what
computation means without committing that meaning to any particular
machine representation. The governing axiom is that meaning is permanent
and representation is disposable. The precise claim is this: SEMANTIC\_ID is the hash of the canonical
semantic IR produced by elaboration, and the hash input includes the
semantic schema version and the normalization rules version:
SEMANTIC\_ID = H(schema\_version || normalization\_version ||
canonical\_graph\_bytes). A stable byte-level graph cannot silently
acquire a new interpretation as the language evolves, because any change
to the schema or the normalization rules changes the identity. Two
surface texts that elaborate into the identical normalized semantic graph
under the same schema and normalization versions share one identity. Programs that elaborate into different
graphs may later be proven equivalent, but that proof is a separate
artifact: source elaborates to canonical semantic IR to SEMANTIC\_ID, and
separately, SEMANTIC\_ID\_A is connected to SEMANTIC\_ID\_B by an
equivalence proof. We do not claim that arbitrary programs can be decided
equivalent; general behavioral equivalence is not computable, and Omega
does not pretend otherwise. What is permanent is the semantic object; what
is disposable is its realization. An optimizer does not change the semantic
object. It realizes the same object into a different target representation,
and carries a refinement or equivalence receipt connecting the realization
to the source semantics.

For the language, we draw four consequences that shape everything below.

First, syntax is not sacred. The concrete text a programmer writes is one
of many possible skins over a meaning. We leave the language free to offer
multiple surfaces, or to change its surface over time, because the surface
was never the program. What the programmer writes is parsed into meaning;
what is stored, checked, imported, and verified is meaning.

Second, equality is semantic, over the canonical graph. Two crumbs whose
surface texts elaborate into the identical normalized semantic graph
share one SEMANTIC\_ID: they are the same semantic object, regardless
of how they were written. They are not necessarily the same crumb: a
crumb is a claim plus its receipts, and different specifications asserted
of the same graph give different CLAIM\_IDs and therefore different
crumbs. We deduplicate semantic objects; we never collapse distinct
claims or distinct evidence. This makes the crumb
store deduplicating by construction: the library cannot fill up with ten
spellings of one idea. Refactoring or optimization that changes the
semantic graph is not identity-preserving; it produces a new semantic
object, connected to the old one by an equivalence receipt where one is
claimed.

Third, verification checks meaning, not text. A receipt does not say ``this
text passed the checks.'' It says ``this meaning was verified against this
specification.'' That is why the crumb module system works at all: when you
import a crumb, you import a meaning plus the evidence for it, not a file.
Files can lie about what they contain. Canonical encoding and hashing do
not prevent that: hashes bind encoded content; they do not establish
semantic correctness. A hash proves the content is the content that was
hashed, nothing more. What hashing gives us is binding: a meaning changed
after hashing no longer matches its hash, so silent substitution is caught.

Fourth, the language and the engine are one idea. The Omega engine's
semantic substrate is the ground the language stands on. The compiler does
not translate text to machine code in the traditional sense; it resolves
meaning and realizes it on the target. The language is the notation; the
semantic substrate is the substance.

In short: most languages ask ``what does this text say?'' Omega asks
``what does this mean?'' and treats everything else as packaging.

\section{Design principles}

\begin{enumerate}
\item \textbf{The verified program is the unit of meaning.} Claims are
admitted through verification, not prose. Semantics defines behavior;
specifications state properties; receipts support those properties under
declared assumptions. The language manipulates meanings, not text;
syntax is a disposable skin over the semantic identity. This is the
Crumb Doctrine, and we make it structural in Omega, not aspirational.

\item \textbf{Design for the writer you have.} We assume the language's code
is mostly written by agents and mostly read by verifiers. We optimize
error messages, tooling, and language rules for machine authors first.
Human readability is a requirement, not the goal.

\item \textbf{Safety is not optional.} We provide no mode in which the safety
guarantees are silently dropped. Any weakening of safety must be explicit,
visible, and auditable. See Section~\ref{sec:escape} for the resolved
answer this raises.

\item \textbf{One idea per mechanism.} Each language feature does one thing.
Features that overlap are a design failure, not richness.

\item \textbf{Name the sacrifice.} Every language gives something up. Go
gave up fine-grained control for simplicity. Rust gave up friendliness for
safety. We state Omega's sacrifice explicitly in
Section~\ref{sec:safety}. We do not leave it to be discovered later by
its users.
\end{enumerate}

\section{The crumb module system}

In Rust, the ecosystem is organized around crates: reusable packages that
programs are built from. In Omega, the corresponding unit is the
\emph{crumb}: a verified program together with its verification receipt. We
define the module and import system of Omega as the crumb store. There is
no separate package registry and no separate proof archive. To import a
crumb is to import its receipt.

Concretely, this means:

\begin{itemize}
\item A crumb is the smallest unit of reuse. Functions, modules, and
libraries are all crumbs: each carries the evidence that it does what it
claims.
\item We permit a program to depend only on verified crumbs. An unverified
program cannot be imported, linked, or built upon. We have the compiler
check this, not policy.
\item A crumb carries several distinct identifiers, and we do not collapse
them. SEMANTIC\_ID is the hash of the canonical meaning together with
the semantic schema version and the normalization rules version
(Section~\ref{sec:semantics}). SPEC\_ID is the
hash of the property or specification being asserted. CLAIM\_ID is the hash
of the semantic ID concatenated with the specification ID:
H(SEMANTIC\_ID || SPEC\_ID). RECEIPT\_ID is the hash of the verifier, the
evidence, the environment, the result, and the CLAIM\_ID the receipt
attests to: the CLAIM\_ID is included in the receipt's hashed payload, so
the receipt is bound to the exact claim it verifies. SEAL\_ID is the hash of the claim ID, the ordered receipt
identifiers, and the store policy ID: H(CLAIM\_ID || ordered\_receipt\_ids
|| policy\_id), an immutable snapshot of the evidence accepted at seal
time. A crumb is the claim plus at least one qualifying receipt: the
logical subject stays stable while evidence accumulates, which matters
when independent laboratories add their own attestations years later. The
same claim may legitimately acquire several receipts: checked by two
independent verifiers, reverified years later, or qualified on two
architectures. Those are new receipts, not new claims. A crumb cannot be
silently substituted or modified without invalidating everything built on
it, because every identifier in the chain is content-addressed.
\item Four rules hold across these identifiers, and we state them together
so the pages agree. The same canonical graph under the same schema and
normalization versions gives the same SEMANTIC\_ID. A different
specification gives a different CLAIM\_ID, because SPEC\_ID is part of
the hash. A different verification event gives a different RECEIPT\_ID
only when event-specific information (timestamp, verifier version, input
digests, environment) is included in the hashed payload: hashing
identical payloads yields identical hashes, and that is the point of the
scheme. A different set of accepted evidence or a different policy gives a
different SEAL\_ID. We deduplicate semantic objects; we never collapse
distinct claims, distinct receipts, or distinct seals.
\item We call the growing lineage of crumbs the \emph{Crumbline}: the
unbroken record of what the machine actually built and verified.
\end{itemize}

Our intent is that the crumb doctrine, which today we must enforce on
agents by external checks, becomes unavoidable. The easy path and the
correct path are the same path, because the language has no other path.

This works because crumbs are semantic objects, not files. A crumb's
identity begins with its SEMANTIC\_ID, the content address of its canonical
meaning. Importing a crumb therefore imports a meaning plus its evidence,
each addressed separately. Two surface texts that elaborate into the
identical normalized semantic graph resolve to the same SEMANTIC\_ID; the
store cannot accumulate ten spellings of one meaning. See
Section~\ref{sec:semantics}.

\section{Safety model}
\label{sec:safety}

We give Omega Rust-grade memory safety: programs cannot corrupt memory,
and the classic classes of memory bugs (use after free, double free, data
races on shared memory, buffer overruns) are ruled out by the language, not
by programmer discipline. We pursue Go-like simplicity in expressing that
safety: the ownership rules are designed to need fewer annotations and
produce clearer errors than Rust's borrow checker, with more inference and
fewer ways to say the same thing.

We state the intended trade up front, per our design principles:
\textbf{we sacrifice low-level control and build speed.} There is no
raw pointer arithmetic in the safe subset, no manual memory layout control,
and no custom allocators. Verification takes time, so builds are slower than
Go's. In exchange, the programmer, human or agent, gets memory safety plus
simplicity plus proof-carrying code. If a use case needs hand-tuned memory
layout or the fastest possible build, Omega is the wrong tool, and we have
the language say so plainly rather than grow a back door.

\section{Concurrency}

Go's most beloved feature is its concurrency model: thousands of tiny,
cheap tasks that are simple to write and simple to reason about. Rust's
concurrency is powerful but widely considered painful. We adopt the
Go-like model as the headline direction: lightweight tasks with simple
syntax, combined with Rust-grade safety guarantees.

We state plainly that this section is direction, not architecture. Before
the concurrency model can carry verification, the following must be
defined: the memory model, the effect model, the ownership and aliasing
model, cancellation semantics, task lifetime rules, channel semantics,
failure propagation semantics, scheduling guarantees, and the treatment
of nondeterminism. The decisive question is not whether the syntax looks
like Go. It is what semantic object a concurrent Omega program denotes,
because verification is meaningless until that object is defined.

\section{Verification as a language feature}

In most languages, verification happens outside the language: test suites,
external analyzers, code review. We make verification a first-class part of
the language, but we do not pretend it is one homogeneous act.
Verification ranges from milliseconds of type checking to hours of
property testing, model checking, differential execution, fuzzing, or
hardware-in-the-loop qualification. One word should not cover all of that.
We therefore divide the pipeline into four stages, and we name them so
their costs stay visible.

\texttt{omega check} establishes well-formedness and declared effects:
the program is coherent and states the effects it claims. This is the
fast, local work: parsing, elaboration, and type checking.

A workspace \texttt{omega build} produces an unsealed realization for
testing: a thing that runs, carrying no admission into the Crumbline.
Building does not verify; it realizes.

\texttt{omega verify} produces claim-specific evidence against the
specification the claim names. The verification profile says what is
checked, how, and at what cost; a receipt for type checking and a receipt
for model checking are different artifacts with different weight.

\texttt{omega seal} admits the artifact and its qualified dependency
closure, minting the crumb: from that point the artifact is addressable
by its identifiers and its closure is on the record. Only a sealed crumb
may enter the crumb store or satisfy a verified dependency.

An unsealed object may exist locally during development; we do not make
the inner edit-build cycle unusable in the name of purity. But we
distinguish the workspace graph from the crumb graph. In the workspace,
sealed and unsealed objects may depend on each other locally: an unsealed
object may depend on unsealed objects during development. The crumb graph
is sealed only. Nothing unsealed may satisfy a sealed dependency, enter
the crumb store, or contribute to a sealed top-level program until the
entire dependency closure meets the required verification policy. That is
how ``no receipt, no crumb'' survives contact with real development.

What is verified is the meaning, not the text. Two surface texts that
elaborate into the identical normalized semantic graph share one
SEMANTIC\_ID, and one verification of that graph serves both. Where the
graphs differ, equivalence must be proven as its own artifact; we do not
declare it by definition.

We draw three consequences from this. First, the Crumbline grows
automatically as a side effect of sealing verified software; no separate step
records what was proven.
Second, the cost of verification is paid once, at seal time, and amortized
over every program that imports the crumb afterward. Third, trust becomes
checkable: any consumer of a crumb can re-examine its receipt without
trusting its author.

We leave the verification machinery itself and the exact contents of a
receipt to the language specification. In this document we establish only
that verification is not optional tooling. It is the grammar.

\section{The foreign boundary}
\label{sec:escape}

We have decided. There is no conventional \texttt{unsafe} inside Omega.

The argument for an escape hatch was practical: hardware interfaces,
foreign code, and extreme performance work may be inexpressible inside the
safe subset. The argument against was doctrinal: a language whose meaning
is the verified program cannot have a region where programs are
unverified but still composable, or the doctrine is a suggestion rather
than a grammar. We resolve this by refusing the premise that the hatch
must live inside the language.

Low-level components, device interfaces, allocators, drivers, SIMD
layouts, DMA, MMIO, and hand-tuned kernels, live outside Omega, in the
languages that can express them. Omega imports them only through sealed
foreign crumbs: boundary artifacts that carry an explicit ABI, a
behavioral contract, platform constraints, an implementation digest, and
qualification evidence. The foreign crumb is sealed by the independent
verifier like any other crumb, against the contract it claims. Omega never
sees the inside; it sees the contract and the receipt, and that is enough
to build upon.

We hold foreign crumbs to an invariant. A foreign crumb must satisfy one
of three conditions: (a) it is formally proven to satisfy the relevant
memory and ABI contract; (b) it executes behind an isolation boundary that
prevents violation of Omega's memory model; or (c) it explicitly weakens
the enclosing crumb's guarantee, and that weakening is visible in the
specification and the receipt. For ordinary foreign components our
preference is isolation plus contracts. Anything looser recreates
\texttt{unsafe} in another file. A behavioral receipt alone cannot
preserve memory safety for foreign code running in the same address space:
without a proof or a boundary, the receipt describes behavior we cannot
enforce, and a crumb built on it makes a guarantee it does not hold.

By ``isolation boundary'' we mean something concrete: the foreign code
cannot read or write Omega-managed memory except through the checked ABI,
and it shares no mutable address space with Omega code. The canonical
mechanism is a hardened OS process. Omega already executes native
candidates this way: the runtime places the code on execute-only pages,
forks a child, runs the candidate there, and kills it on timeout
(\texttt{src/runtime/rx\_omega.c}). That is crash and hang containment,
not yet a security sandbox: the current child has no seccomp filter, no
namespace or filesystem isolation, no resource limits, and no privilege
dropping. The hardened form closes those gaps. The child starts with no
ambient file descriptors, no network, a seccomp allowlist, namespace and
filesystem isolation where the platform provides it, explicit memory and
CPU limits, no-new-privileges, and a killable timeout. We keep the
separate-process model rather than WASM because Omega must test the
actual native realization it may promote; validating WASM is not
validating the AArch64 code that will run. WASM remains a separate
execution class for portable tools. The architectural rule is: foreign
code is data until a disposable, capability-empty native worker verifies
it; only receipt-bound verified code may cross into the resident
execution path.

One caution about the mechanism itself. A forked child inherits the
readable contents of the parent's memory; closing descriptors and adding
syscall restrictions does not erase that information. Worker
initialization must therefore start from a clean, minimal state: the
child is given only what the task requires, and the rest of the parent's
address space is cleared or dropped before the foreign code runs.

A promotion rule follows from the invariant. Foreign code qualified
through isolation remains isolated in production: if the receipt that
admits a component rests on an isolation boundary, the component stays
behind that boundary in every deployment. Resident-address-space
execution requires the relevant memory/ABI proof, or an explicitly
propagated weakening of the enclosing guarantee under condition (c)
above. We do not let a component ride an isolation receipt into the
resident address space; that would be \texttt{unsafe} wearing a sealed
label.

A weakening under (c) is not a note in prose. It is a requirement on
CLAIM\_ID and spec composition that the weakened guarantee is part of
the specification itself: SPEC\_ID is computed over the weakened
specification, so every dependent CLAIM\_ID is derived from the weakened
claim, and the admission check rejects any seal whose dependency closure
treats the component as unweakened. The weakening propagates mechanically
through the identifier chain and the admission checks, not by convention.

This preserves the doctrine better than embedding an unverified language
inside the verified one. There is no region of Omega where verification
stops applying, because the unverified work never enters Omega at all. It
stands at the boundary, labeled, receipted, and accountable.

\section{Bootstrapping path}

The foundation crumbs, the verified core the language is built from, must
be verified the hard way first: written carefully, checked thoroughly, with
no existing Crumbline to stand on. This is slow. There is no way to skip it,
because a crumb store is only as trustworthy as its first entries.

The expected shape is therefore one we accept: a slow, careful beginning,
then compounding returns. Each verified foundation crumb makes the next
layer cheaper to verify, because new crumbs are built from old verified
ones.
The library grows slowly at first and then quickly, and every entry pays
for itself across all future work. We accept patience at the start as part
of the design.

\section{Objections}

We have stated our proposal. We now consider the objections we expect to
be raised against it, and answer each in turn.

\subsection*{Objection 1: This is only Rust with nicer syntax.}

We answer that this description would be a smaller idea, and it misses the
proposal. A nicer Rust would still manipulate text and hope meaning
follows. Our distinguishing claim is different in kind: verification as a
four-stage pipeline (\texttt{check}, \texttt{build}, \texttt{verify},
\texttt{seal}), the crumb as the unit of reuse, and meaning rather than
text as the object the language manipulates. We are not improving Rust's
answers. We are asking a different question.

\subsection*{Objection 2: Or Go with a borrow checker.}

The same answer applies. Go's simplicity is worth borrowing, and we borrow
it openly. But a Go with stricter checking would still be a language about
text, with verification as an afterthought. We make verification and sealing
structural requirements of reusable composition, and we make the module
system the store of verified meanings. The resemblance to Go is in the feel of writing concurrent code,
not in the nature of the language.

\subsection*{Objection 3: Most tasks do not need any of this.}

We agree. Omega is not a general-purpose scripting language. Where a task
does not need safety guarantees or verification, we recommend the simpler
tool. A language that insists on proving everything is the wrong instrument
for work that needs proving nothing. We do not propose Omega for all
programming. We propose it for the programming that has to be right.

\subsection*{Objection 4: A language about text would be simpler and just as good.}

It would be simpler to build, and it would not be as good. Text is
packaging: two surface texts that elaborate into the identical normalized
semantic graph share one identity, and a language that cannot see this
will verify the same meaning twice, store ten spellings of one idea, and
mistake reformatting for change. We do not claim to decide equivalence of
arbitrary programs; where the semantic graphs differ, equivalence is a
separate proven artifact. We accept the cost of canonicalizing meaning
because the alternative is a system that confuses the wrapper for the
thing wrapped.

\subsection*{Objection 5: This is academic theorizing, not engineering.}

We measure Omega's success in systems shipped, not papers published. It is
a working language for building real systems, or it is nothing. The
proposal is written down precisely so that it can be tested against
reality, and we will discard whatever reality rejects.

\subsection*{Objection 6: Why should anyone believe these claims before the language exists?}

They should not. This document is a design proposal dated October 2026,
and every claim in it is provisional until we have built, measured, and
ratified it. We propose; we do not assert. The objections above are
welcome, and we expect more. A proposal that cannot survive its critics
does not deserve to be built.

\subsection*{Objection 7: This has all been done before.}

Much of it has, and we say so plainly. Unison addressed code by the hash
of its content. Dafny, F*, SPARK, and Lean made verification a
first-class concern of the language rather than an afterthought.
Proof-carrying code, introduced by Necula in 1997, is the direct ancestor
of our maxim that the receipt travels with the code. We do not claim to
be the first to put verification inside a programming language.

What we propose is the integration of four things these systems do not
combine: semantic identity as the unit of reuse, mandatory
evidence-bearing dependencies, an admission policy that decides what may
be built upon, and a provenance record of everything the machine has built
and verified. Each of those has precedent. Their combination, made
structural rather than optional, is the proposal.

\end{document}
